\documentclass[../main.tex]{subfiles}
\begin{document}
\section{Truth-uniform query error contracts}
\label{sec:query_contracts}
Tuple marginals alone do not specify a query's error. Source identity is
needed to construct its lineage, and the exact truth profile determines
which sources can err in each direction. We now ask for a contract valid
over every exact truth profile giving the specified query answer. This
differs from exact probability evaluation for one fixed profile, though it
uses the same Boolean lineage.

Fix a nonconstant monotone function $F:\{0,1\}^n\to\{0,1\}$ on distinct
sources. In the one-sided false-positive model,
$Y_i^+(x)=x_i\lor E_i$ with independent $E_i\sim\mathrm{Bernoulli}(\alpha_i)$.
In the separate one-sided false-negative model,
$Y_i^-(x)=x_i\land(1-E_i')$ with independent
$E_i'\sim\mathrm{Bernoulli}(\beta_i)$. Define
\[
 \mathcal E_0(F;\boldsymbol\alpha)
 =\max_{x:F(x)=0}\Pr\{F(Y^+(x))=1\},\qquad
 \mathcal E_1(F;\boldsymbol\beta)
 =\max_{x:F(x)=1}\Pr\{F(Y^-(x))=0\}.
\]
No distribution over exact truth is assumed. Every source assignment is
admissible here. This holds for class-induced equality constraints after
quotienting; additional constraints such as transitivity may exclude the
maximizers, leaving upper bounds rather than sharp contracts. For constant
functions the absent conditioning class has no error rate; setting its
extremal envelope to zero is only a mathematical convention.

\subsection{Boundary certificates give sharp one-sided contracts}
Let $\mathcal B_0$ be the inclusion-maximal false vectors and
$\mathcal B_1$ the inclusion-minimal true vectors. The one-sided sensitivities are
\[
 s_b(F)=\max_{x:F(x)=b}|\{i:F(x\mathbin\oplus e_i)\ne F(x)\}|,
 \qquad b\in\{0,1\}.
\]
For monotone functions these equal the corresponding certificate
complexities, a standard fact~\cite{blanc2022counterfactuals}. We use it to
express exact one-sided error envelopes; sensitivity-based bounds for
two-sided noise are also established~\cite[Lemma~15]{gopalan2016smooth}.

\begin{theorem}[Independent-source boundary contract]
\label{thm:query_contract}
For the one-sided models above,
\begin{align}
 \mathcal E_0(F;\boldsymbol\alpha)
 &=\max_{b\in\mathcal B_0}
       \left[1-\prod_{i:b_i=0}(1-\alpha_i)\right],\label{eq:hetero_contract_fp}\\
 \mathcal E_1(F;\boldsymbol\beta)
 &=\max_{b\in\mathcal B_1}
       \left[1-\prod_{i:b_i=1}(1-\beta_i)\right].\label{eq:hetero_contract_fn}
\end{align}
For homogeneous rates these specialize to
\begin{equation}
 \boxed{\mathcal E_0(F;\alpha)=1-(1-\alpha)^{s_0(F)}},\qquad
 \boxed{\mathcal E_1(F;\beta)=1-(1-\beta)^{s_1(F)}}.
 \label{eq:homogeneous_contract}
\end{equation}
\end{theorem}
\begin{proof}
Couple all truth profiles using the same independent uniforms for the
source errors. If $x\leq x'$, then $Y^+(x)\leq Y^+(x')$ pointwise.
Every false $x$ extends to a maximal false $b$, so monotonicity reduces
the maximum false-positive probability to $\mathcal B_0$. At such a $b$,
flipping any remaining zero makes $F$ true. The query stays false exactly
when all those zeros stay zero, giving \eqref{eq:hetero_contract_fp}.
Shrinking true vectors to minimal true vectors gives the dual formula.

Every zero of a maximal false vector is pivotal. Conversely, a pivotal
zero at an arbitrary false vector remains zero and pivotal at every maximal
false extension: setting it to one would already force the query true.
Thus $s_0$ is the maximum number of zeros on $\mathcal B_0$. The dual argument
gives $s_1$ as the maximum number of ones on $\mathcal B_1$, proving
\eqref{eq:homogeneous_contract}.
\end{proof}

These are elementary consequences of monotone boundary certificates, not
a new general inference algorithm. In a two-sided independent channel the
same expressions are upper bounds: suppressing false negatives can only
increase false positives, and conversely. They need not be equalities.
For $F=X_1\land X_2$ with $\alpha\leq1-\beta$, the worst two-sided FPR is
$\alpha(1-\beta)$, whereas the one-sided formula gives $\alpha$.

\begin{theorem}[Marginal-only dependence bound]
\label{thm:dependent_contract}
Without independence, but with per-source marginal FP bounds $\alpha_i$,
the sharp envelope over exact false profiles and all couplings is
\[
 \max_{b\in\mathcal B_0}\min\left\{1,\sum_{i:b_i=0}\alpha_i\right\}.
\]
For homogeneous bounds it is $\min\{1,\alpha s_0(F)\}$. The FN statement
is dual, using $\mathcal B_1$, $\beta_i$, and $s_1(F)$.
\end{theorem}
\begin{proof}
Extend a false input to a maximal false $b$. Its remaining zeros form
a zero certificate already present in the input. Error requires at least
one of these coordinates to flip, so the union bound gives the result.
At a maximizing boundary profile every one of those flips suffices. Events
of prescribed probabilities can attain union probability equal to their
sum capped at one: use disjoint intervals when the sum is at most one,
and consecutive intervals on a unit circle otherwise. This coupling
attains the bound. The true-input argument is dual.
\end{proof}
The upper bounds also hold with two-sided noise by the same certificate
argument. Sharpness here is over allowed couplings, not a statement about
any particular shared hash construction.

\begin{corollary}[A shared random rate]
\label{cor:shared_rate}
Fix an exact false input. Suppose that, conditional on a construction
statistic $C$, the distinct absent sources are independent with common FPR
$P(C)$, and true sources are always reported. Then
\[
 \Pr\{F(Y)=1\}
 \leq\mathbb E[1-(1-P(C))^{s_0(F)}]
 \leq1-(1-\mathbb E[P(C)])^{s_0(F)}.
\]
\end{corollary}
\begin{proof}
Apply \cref{thm:query_contract} conditionally and average. The function
$p\mapsto1-(1-p)^s$ is concave for integer $s\geq1$, so Jensen's
inequality gives the second bound.
\end{proof}
This conditional model permits marginal dependence. For example two absent
sources can have conjunction probability $\mathbb E[P^2]$, exceeding
$(\mathbb E[P])^2$, while the truth-uniform bound still holds. A uniform
mean-rate bound over exact inputs therefore yields a uniform contract;
equality requires additional assumptions. Arbitrary dependence cannot be
replaced by this particular common-rate model.

\subsection{Binary witnesses of occurrence at most two}
Write lineage in normalized positive disjunctive normal form
$F=\bigvee_j\bigwedge_{i\in S_j}X_i$: identify aliases, deduplicate sources
within each witness, remove duplicate witnesses, and absorb each witness
that contains another. Then $s_1=\max_j|S_j|$, and $s_0$ is the largest
size of an \emph{inclusion-minimal} hitting set of the witness family.
It is not the minimum hitting-set size.

For binary witnesses, let $G$ be the simple graph whose edges are the
supports $S_j$. False vectors correspond to independent sets; maximal
false vectors correspond to maximal independent sets. Consequently
\begin{equation}
 s_0(F_G)=|V(G)|-i(G),
 \label{eq:independent_domination}
\end{equation}
where $i(G)$ is the minimum cardinality of a maximal independent set, also
called the independent domination number. Computing it is hard already
on restricted cubic graphs~\cite{liu2015independent}; merely limiting
witness size to two does not make contracts easy.

\begin{theorem}[A tractable query class]
\label{thm:degree_two_contract}
If normalized witness width and source occurrence are at most two, let $U_1$
be the singleton witnesses and $G_c$ the remaining path or cycle components,
of orders $n_c$. Then
\[
 s_0=|U_1|+\sum_c\left(n_c-\left\lceil n_c/3\right\rceil\right),
 \qquad s_1=\max_j|S_j|.
\]
After normalization, these values require a linear number of graph operations.
\end{theorem}
\begin{proof}
The occurrence restriction makes graph degree at most two. An independent
dominating vertex covers at most three vertices, giving
$i(G_c)\geq\lceil n_c/3\rceil$. For a path, take every third vertex starting
at the second when $n_c\equiv0,2\pmod3$, and at the first when
$n_c\equiv1\pmod3$. For a cycle, distribute $r=\lceil n_c/3\rceil$ selected
vertices with cyclic gaps two or three summing to $n_c$; this is possible
because $2r\leq n_c\leq3r$. These sets are independent and dominating.
Equation~\eqref{eq:independent_domination} adds over components. Every
singleton witness must be hit, and normalization removes its incident
supersets, giving the singleton contribution.
\end{proof}

For a target query FPR $\delta\in(0,1)$, homogeneous independent one-sided
source rates satisfying
\begin{equation}
 \alpha\leq1-(1-\delta)^{1/s_0}
 \label{eq:calibrated_alpha}
\end{equation}
suffice when $s_0>0$. Under arbitrary dependence, $\alpha\leq\delta/s_0$
suffices instead. A 12-cycle has $s_0=8$ and $s_1=2$: a per-source FPR
of $0.05$ permits worst-case query FPR $1-0.95^8\approx0.337$, whereas
the all-false profile has a much smaller error requiring two adjacent
false positives. The extremal false profile has true sources at
$\{0,3,6,9\}$ and every other source false. Evaluating only the all-false
profile therefore does not validate a truth-uniform query contract.

For comparison, a fixed-profile query probability on a path has a standard
two-state recurrence. If $a,b$ are the probabilities of no adjacent ones
and a final zero or one, appending a source with success probability $p$
gives $a'=(a+b)(1-p)$ and $b'=ap$. Condition on the first bit to handle a
cycle. This uses linear arithmetic work and constant auxiliary state;
the tractability is established finite-state inference, not a new algorithm.
Heterogeneous contracts can likewise use standard weighted independent
domination dynamic programming, rather than the homogeneous closed form.

\subsection{Why preserving tuple rates does not preserve inference cost}
Canonical symmetry can turn read-once lineage into shared-source lineage.
For a simple cubic graph $G$, let the exact relation $D(v,a,b,c)$ list each
vertex's three neighbors. Consider the fixed positive query
\[
 \exists v,a,b,c:\ D(v,a,b,c)\land R(v,a)\land R(v,b)\land R(v,c).
\]
Suppose all queried $R$ tuples are exact nonmembers with FPR $p$.
Independent directed sources yield query FPR $1-(1-p^3)^{|V|}$.
After symmetry identifies opposite orientations, each undirected edge is
one independent bit used twice. Each tuple still has FPR $p$ and each
witness still has probability $p^3$. Yet the query is false exactly when
the zero edges form an edge cover. At $p=1/2$ its FPR is
\[
 1-\frac{\#\operatorname{EdgeCovers}(G)}{2^{|E|}}.
\]
Edge-cover counting on simple planar cubic graphs is already
\#P-hard~\cite{cai2012holographic,amarilli2022matchings}.
This is an application of known counting hardness, not a new dichotomy.
It explains why preserving a structural invariant at unchanged tuple rates
need not preserve the cost of exact query inference. A compact worst-case
contract and a fixed-profile exact probability are different objectives.
\end{document}
